% ============================================================================
%  Laboratorio di Programmazione -- Lezione 1 (15 settembre 2026)
%  Compilare con:  pdflatex -shell-escape Lezione01_Benvenuti_in_Java.tex
%  Richiede: pacchetto `mermaid` (CTAN: ltmermaid) + Mermaid CLI `mmdc` nel PATH
%  NB: ogni frame che contiene codice (lstlisting) o un diagramma (mermaid)
%      deve essere dichiarato [fragile].
% ============================================================================
\documentclass[10pt,aspectratio=169]{beamer}
\usetheme{Madrid}
\usecolortheme{default}

% `provide=*` carica la localizzazione italiana dai file .ini di babel: funziona anche
% su installazioni TeX prive di texlive-lang-italian; con una TeX Live completa e`
% equivalente al classico \usepackage[italian]{babel}.
\usepackage[italian,provide=*]{babel}
\usepackage[utf8]{inputenc}

% Localizzazione dei nomi dei blocchi Beamer
\uselanguage{italian}
\languagepath{italian}
\deftranslation[to=italian]{Definition}{Definizione}
\deftranslation[to=italian]{Theorem}{Teorema}
\deftranslation[to=italian]{Example}{Esempio}
\deftranslation[to=italian]{Proof}{Dimostrazione}
\deftranslation[to=italian]{Corollary}{Corollario}
\deftranslation[to=italian]{Lemma}{Lemma}

\usepackage{listings}
\usepackage{xcolor}
\usepackage{hyperref}
\usepackage{array}
\usepackage{booktabs}
\usepackage{mermaid}

% Mermaid: sfondo trasparente e dimensione massima adatta a una slide 16:9
\MermaidRendererOptions{-b transparent}
\MermaidAdjustBoxOpts{max width=0.95\linewidth,max height=0.62\textheight,center}

% Configurazione colori per il codice
\definecolor{codegreen}{rgb}{0,0.6,0}
\definecolor{codegray}{rgb}{0.5,0.5,0.5}
\definecolor{codepurple}{rgb}{0.58,0,0.82}
\definecolor{backcolour}{rgb}{0.95,0.95,0.92}
\definecolor{keywordblue}{rgb}{0.0,0.0,0.7}

% Stile per snippet Java
\lstdefinestyle{javastyle}{
    backgroundcolor=\color{backcolour},
    commentstyle=\color{codegreen}\itshape,
    keywordstyle=\color{keywordblue}\bfseries,
    numberstyle=\scriptsize\color{codegray},
    stringstyle=\color{codepurple},
    basicstyle=\ttfamily\footnotesize,
    breakatwhitespace=false,
    breaklines=true,
    captionpos=b,
    keepspaces=true,
    numbers=none,
    numbersep=5pt,
    showspaces=false,
    showstringspaces=false,
    showtabs=false,
    tabsize=2,
    language=Java,
    literate={à}{{\`a}}1 {è}{{\`e}}1 {é}{{\'e}}1 {ì}{{\`i}}1
             {ò}{{\`o}}1 {ù}{{\`u}}1 {È}{{\`E}}1
}
\lstset{style=javastyle}

% Stile per righe di terminale
\lstdefinestyle{shell}{
    language=bash,
    backgroundcolor=\color{black!85},
    basicstyle=\ttfamily\footnotesize\color{white},
    keywordstyle=\color{white},
    stringstyle=\color{white},
    commentstyle=\color{green!70},
    numbers=none,
    frame=none,
    showstringspaces=false
}

\title[Lezione 1 -- Benvenuti in Java]{Lezione 1 -- Benvenuti in Java}
\subtitle{Il corso, la piattaforma Java, il primo programma}
\author[Mattia Fonisto]{Mattia Fonisto}
\institute[Lab.\ di Programmazione]{Laboratorio di Programmazione -- Ingegneria Informatica, II anno\\
           Università degli Studi di Napoli Federico II -- Sede di San Giovanni a Teduccio}
\date{15 settembre 2026}

\begin{document}

\frame{\titlepage}

% ========================================================================
\begin{frame}{Agenda della lezione (2h)}
\begin{enumerate}
    \item \textbf{Il corso} \hfill (\emph{15 min})
    \begin{itemize}
        \item obiettivi, percorso, organizzazione, homework
    \end{itemize}
    \item \textbf{Java: perché e come} \hfill (\emph{25 min})
    \begin{itemize}
        \item ``scrivi una volta, esegui ovunque'': bytecode e JVM
        \item JDK, JRE, versioni; che cosa installare
    \end{itemize}
    \item \textbf{Il primo programma} \hfill (\emph{30 min})
    \begin{itemize}
        \item anatomia di \texttt{HelloWorld.java}
        \item edit \textrightarrow{} \texttt{javac} \textrightarrow{} \texttt{java}: dal vivo, da terminale
        \item gli errori che farete (tutti) la prima settimana
    \end{itemize}
    \item \textbf{Esercitazione guidata} -- programmi interattivi \hfill (\emph{40 min})
    \begin{itemize}
        \item \texttt{System.out.printf}, \texttt{Scanner}, metodi \texttt{static}
        \item \texttt{SommaNumeri} commentato insieme, poi \texttt{Media} scritto insieme, passo passo
    \end{itemize}
    \item \textbf{Wrap-up e homework} \hfill (\emph{10 min})
\end{enumerate}
\end{frame}

% ========================================================================
\section{Il corso}

\begin{frame}{Obiettivi del corso}
Al termine del corso saprete:
\begin{itemize}
    \item \textbf{progettare} piccoli sistemi software in modo modulare, ragionando per \emph{astrazioni};
    \item \textbf{programmare a oggetti} in Java: classi, incapsulamento, ereditarietà, polimorfismo, interfacce;
    \item usare alcune \textbf{tecniche avanzate}: ricorsione, eccezioni, tipi generici, API standard;
    \item implementare le \textbf{strutture dati fondamentali}: pila, coda, liste collegate;
    \item leggere e scrivere \textbf{file di testo}.
\end{itemize}

\vspace{0.3cm}
\begin{block}{È un laboratorio}
Ogni lezione contiene codice scritto \emph{insieme} e codice scritto \emph{da voi}. Portate il portatile con il JDK installato (oggi vediamo come).
\end{block}
\end{frame}

% ========================================================================
\begin{frame}[fragile]{Il percorso: 13 lezioni, 4 blocchi}
\begin{mermaid}
flowchart LR
  A["<b>A · Fondamenti</b><br/>L1–L4<br/>Java e JVM · tipi · String<br/>metodi · array · modularità"]
  B["<b>B · Oggetti</b><br/>L5–L9<br/>classi · this/static · immutabilità<br/>ereditarietà · polimorfismo<br/>Object · interfacce · package"]
  C["<b>C · Avanzato</b><br/>L10–L12<br/>eccezioni · ricorsione<br/>ricerca binaria · generici<br/>Comparable · pile, code, liste"]
  D["<b>D · Chiusura</b><br/>L13<br/>file di testo<br/>ripasso e simulazione"]
  A --> B --> C --> D
  style A fill:#dbeafe,stroke:#1d4ed8
  style B fill:#dcfce7,stroke:#15803d
  style C fill:#fef3c7,stroke:#b45309
  style D fill:#fce7f3,stroke:#be185d
\end{mermaid}

\vspace{0.2cm}
\small Martedì, 2h, dal 15 settembre al 15 dicembre (l'8 dicembre è festivo). Il programma completo, con date e contenuti di ogni lezione, sarà pubblicato sul team Microsoft Teams del corso (seguirà comunicazione).
\end{frame}

% ========================================================================
\begin{frame}{Come funziona una lezione -- e il corso}
\begin{columns}[T]
\column{0.5\textwidth}
\textbf{Struttura tipica (2h)}
\begin{itemize}\itemsep=2pt
    \item 10--15' ripasso e correzione homework
    \item 35--45' concetti nuovi, con esempi \emph{eseguiti}
    \item 50--60' esercitazione
    \item 5--10' wrap-up e homework
\end{itemize}

\vspace{0.3cm}
\textbf{Strumenti}
\begin{itemize}\itemsep=2pt
    \item JDK 25 (LTS, Eclipse Temurin)
    \item Visual Studio Code con l'\emph{Extension Pack for Java} (lo installiamo oggi)
    \item terminale per capire cosa fa l'IDE ``sotto''
\end{itemize}

\column{0.5\textwidth}
\textbf{Homework}
\begin{itemize}\itemsep=2pt
    \item uno a settimana, da svolgere prima della lezione successiva
    \item non obbligatori ma \textbf{fortemente consigliati}
    \item all'inizio di ogni lezione chiederò a qualcuno di \textbf{presentare la propria soluzione}: la discutiamo insieme
\end{itemize}

\vspace{0.3cm}
\textbf{Testo di riferimento}
\begin{itemize}\itemsep=2pt
    \item Deitel \& Deitel, \emph{Programmare in Java}, 11\textsuperscript{a} ed., Pearson
    \item slide, codice ed esercizi sul team \textbf{Microsoft Teams} del corso (seguirà comunicazione)
\end{itemize}
\end{columns}
\end{frame}

% ========================================================================
\section{Java: perché e come}

\begin{frame}{Che cos'è Java}
\begin{itemize}
    \item Nasce in Sun Microsystems (progetto \emph{Green}, 1991; ``Oak''), annunciato ufficialmente il 23 maggio 1995. Dal 2006--2007 l'implementazione di riferimento è \emph{open source} (OpenJDK).
    \item È un linguaggio \textbf{orientato agli oggetti}, \textbf{fortemente tipizzato}, con \textbf{gestione automatica della memoria} (garbage collector).
    \item È definito da un documento, la \emph{Java Language Specification} (JLS); la piattaforma è \textbf{Java SE} (Standard Edition), un insieme di specifiche mantenute dal Java Community Process.
\end{itemize}

\vspace{0.2cm}
\begin{block}{La filosofia: \emph{Write Once, Run Anywhere}}
Un programma Java viene scritto e compilato \textbf{una volta} e poi eseguito \textbf{ovunque} esista una Java Virtual Machine: Windows, macOS, Linux, Android, server, schede embedded\ldots
\end{block}

\vspace{0.2cm}
\small Come è possibile? La risposta è nella parola \textbf{bytecode}.
\end{frame}

% ========================================================================
\begin{frame}[fragile]{Compilazione tradizionale (es.\ C/C++) vs Java}
\MermaidAdjustBoxOpts{max width=0.95\linewidth,max height=0.46\textheight,center}
\begin{columns}[T]
\column{0.5\textwidth}
\centering\textbf{C++: un compilatore (e un eseguibile) per ogni architettura}
\begin{mermaid}
flowchart LR
  S1["sorgente<br/>.cpp"] --> C1["compilatore<br/>x86"] --> O1["eseguibile<br/>x86"]
  S1 --> C2["compilatore<br/>ARM"] --> O2["eseguibile<br/>ARM"]
  S1 --> C3["compilatore<br/>RISC-V"] --> O3["eseguibile<br/>RISC-V"]
\end{mermaid}

\column{0.5\textwidth}
\centering\textbf{Java: un solo bytecode, una JVM per architettura}
\begin{mermaid}
flowchart LR
  S2["sorgente<br/>.java"] -->|javac| B["bytecode<br/>.class"]
  B --> J1["JVM x86"]
  B --> J2["JVM ARM"]
  B --> J3["JVM RISC-V"]
  style B fill:#fef3c7,stroke:#b45309
\end{mermaid}
\end{columns}
\MermaidAdjustBoxOpts{max width=0.95\linewidth,max height=0.62\textheight,center}

\vspace{0.3cm}
\begin{block}{Il punto}
\texttt{javac} \textbf{non} genera codice macchina ma \emph{bytecode}, per un processore virtuale. Si compila \textbf{una volta}; a tradurlo è la JVM installata su ogni macchina.
\end{block}
\end{frame}

% ========================================================================
\begin{frame}[fragile]{Il ciclo di vita di un programma Java: le 5 fasi}
\begin{mermaid}
flowchart LR
  E["<b>1. Edit</b><br/>editor / IDE<br/>Prova.java"]
  C["<b>2. Compile</b><br/>javac Prova.java<br/>→ Prova.class"]
  L["<b>3. Load</b><br/>Class Loader<br/>carica il .class in RAM"]
  V["<b>4. Verify</b><br/>Bytecode Verifier<br/>controlli di sicurezza"]
  X["<b>5. Execute</b><br/>Execution Engine<br/>interprete + JIT"]
  E --> C --> L --> V --> X
  subgraph JVM ["java Prova  (Java Virtual Machine)"]
    L
    V
    X
  end
\end{mermaid}

\vspace{0.3cm}
\begin{columns}[T]
\column{0.5\textwidth}
\textbf{A tempo di compilazione}
\begin{itemize}\itemsep=0pt
    \item controllo sintattico e dei \emph{tipi}
    \item un file \texttt{.class} \textbf{per ogni classe}
\end{itemize}
\column{0.5\textwidth}
\textbf{A tempo di esecuzione}
\begin{itemize}\itemsep=0pt
    \item il \emph{collegamento} è dinamico: le classi sono caricate solo quando servono
    \item ogni classe caricata viene verificata
\end{itemize}
\end{columns}
\end{frame}

% ========================================================================
\begin{frame}[fragile]{Dentro la JVM}
\MermaidAdjustBoxOpts{max width=0.95\linewidth,max height=0.46\textheight,center}
\begin{mermaid}
flowchart LR
  F["file .class<br/>(bytecode)"]
  CL["<b>Class Loader</b><br/>carica e collega le classi<br/>solo quando servono"]
  BV["<b>Bytecode Verifier</b><br/>bytecode conforme?<br/>nessun cast illegale,<br/>nessun overflow di stack"]
  subgraph EE ["Execution Engine"]
    I["<b>Interprete</b><br/>istruzione per istruzione"]
    JIT["<b>JIT Compiler</b><br/>compila i blocchi 'caldi'<br/>in codice macchina"]
  end
  F --> CL --> BV --> EE
  I -.->|codice eseguito spesso| JIT
  style F fill:#fef3c7,stroke:#b45309
\end{mermaid}
\MermaidAdjustBoxOpts{max width=0.95\linewidth,max height=0.62\textheight,center}

\vspace{0.1cm}
\footnotesize
\begin{columns}[T]
\column{0.5\textwidth}
\begin{itemize}\itemsep=2pt
    \item La JVM è \textbf{processore virtuale + sistema operativo virtuale}: un ambiente completo e isolato.
    \item Il \emph{verifier} impedisce al codice di violare l'integrità del sistema (es.\ trasformare un intero in un puntatore).
\end{itemize}
\column{0.5\textwidth}
\begin{itemize}\itemsep=2pt
    \item L'\emph{interprete} parte subito; il \emph{JIT} ottimizza le parti eseguite più spesso: oggi Java è competitivo con i linguaggi compilati.
    \item La \textbf{memoria è gestita dalla JVM}: creiamo oggetti con \texttt{new}, il \emph{garbage collector} li libera.
\end{itemize}
\end{columns}
\end{frame}

% ========================================================================
\begin{frame}{JDK, JRE, versioni: che cosa installare}
\begin{columns}[T]
\column{0.5\textwidth}
\textbf{JDK -- Java Development Kit}
\begin{itemize}\itemsep=2pt
    \item tutto ciò che serve per \textbf{sviluppare}: compilatore \texttt{javac}, JVM \texttt{java}, strumenti (\texttt{jshell}, \texttt{javadoc}, \texttt{jar}\ldots) e la libreria standard
\end{itemize}

\vspace{0.2cm}
\textbf{JRE -- Java Runtime Environment}
\begin{itemize}\itemsep=2pt
    \item solo la JVM e la libreria: basta per \textbf{eseguire}
    \item dalla versione 11 non si distribuisce più separatamente: \textbf{installate il JDK}
\end{itemize}

\vspace{0.2cm}
\textbf{La libreria standard}
\begin{itemize}\itemsep=2pt
    \item migliaia di classi organizzate in \emph{package}: \texttt{java.lang} (String, Math, System\ldots), \texttt{java.util} (Scanner, collezioni), \texttt{java.io} (file)\ldots
    \item documentazione ufficiale: \emph{docs.oracle.com} \textrightarrow{} Java SE 25 API
\end{itemize}

\column{0.5\textwidth}
\textbf{Versioni}
\begin{itemize}\itemsep=2pt
    \item una release ogni 6 mesi (marzo e settembre): due numeri di versione all'anno
    \item \textbf{LTS} (Long Term Support) ogni 2 anni, cioè una versione ogni quattro: 17 (2021), 21 (2023), \textbf{25} (2025)
    \item per il corso usiamo tutti \textbf{JDK 25}: stessa versione, stessi comportamenti, stessi messaggi d'errore -- distribuzione consigliata: Eclipse Temurin (Adoptium)
\end{itemize}
\end{columns}
\end{frame}

% ========================================================================
\begin{frame}[fragile]{Verifichiamo l'installazione (adesso, tutti)}
Da terminale (macOS/Linux) o prompt dei comandi / PowerShell (Windows):
\begin{lstlisting}[style=shell]
$ java -version
openjdk version "25.0.2" 2026-01-20 LTS
OpenJDK Runtime Environment Temurin-25.0.2+10 (build 25.0.2+10-LTS)
OpenJDK 64-Bit Server VM Temurin-25.0.2+10 (build 25.0.2+10-LTS, mixed mode)

$ javac -version
javac 25.0.2
\end{lstlisting}

\begin{columns}[T]
\column{0.5\textwidth}
\textbf{Se \texttt{java} funziona ma \texttt{javac} no}
\begin{itemize}\itemsep=0pt
    \item avete un JRE, non un JDK: installate il JDK
\end{itemize}
\textbf{Se nessuno dei due funziona}
\begin{itemize}\itemsep=0pt
    \item JDK non installato, oppure non è nel \texttt{PATH}
    \item su Windows verificate anche la variabile \texttt{JAVA\_HOME}
\end{itemize}
\column{0.5\textwidth}
\textbf{Download del JDK 25}
\begin{itemize}\itemsep=0pt
    \item \url{https://adoptium.net} \textrightarrow{} Temurin 25 (LTS)
    \item Windows: scegliete l'installer \texttt{.msi} e spuntate \emph{Set JAVA\_HOME} e \emph{Add to PATH}
    \item macOS: \texttt{.pkg}; Linux: pacchetto della distribuzione (\texttt{openjdk-25-jdk})
\end{itemize}
\textbf{IDE}
\begin{itemize}\itemsep=0pt
    \item useremo sempre \textbf{Visual Studio Code}: prossima slide
\end{itemize}
\end{columns}
\end{frame}

% ========================================================================
\begin{frame}[fragile]{Visual Studio Code per Java: installazione in 4 passi}
\small
\begin{enumerate}\itemsep=2pt
    \item \textbf{Installare VS Code} da \url{https://code.visualstudio.com}
    \item \textbf{Installare l'estensione}: icona \emph{Extensions} (\texttt{Ctrl+Shift+X}), cercare \textbf{Extension Pack for Java} (Microsoft), \emph{Install}. Porta con sé supporto al linguaggio, debugger, test runner e gestione progetti.
    \item \textbf{Collegare il JDK}: \texttt{Ctrl+Shift+P} \textrightarrow{} \emph{Java: Configure Java Runtime}; deve comparire il JDK 25. Se non lo trova, impostare in \emph{Settings} \texttt{java.jdt.ls.java.home} sulla cartella del JDK.
    \item \textbf{Aprire una cartella} (\emph{File \textrightarrow{} Open Folder}), non un singolo file: VS Code compila da solo ogni \texttt{.java} che contiene.
\end{enumerate}

\MermaidAdjustBoxOpts{max width=0.9\linewidth,max height=0.2\textheight,center}
\begin{mermaid}
flowchart LR
  A["1. VS Code"] --> B["2. Extension Pack<br/>for Java"] --> C["3. Configure Java<br/>Runtime → JDK 25"] --> D["4. Open Folder<br/>con i sorgenti"] --> E["▶ Run · F5 Debug"]
  style E fill:#dcfce7,stroke:#15803d
\end{mermaid}
\MermaidAdjustBoxOpts{max width=0.95\linewidth,max height=0.62\textheight,center}

\begin{block}{Come si esegue un programma}
Sopra il \texttt{main} compaiono i link \textbf{Run | Debug} (oppure \texttt{F5}). L'output appare nel pannello \emph{Terminal}: sotto, VS Code lancia lo stesso \texttt{javac} + \texttt{java} che oggi vediamo a mano.
\end{block}
\end{frame}

% ========================================================================
\section{Il primo programma}

\begin{frame}[fragile]{\texttt{HelloWorld.java}}
\begin{lstlisting}
/* Prima classe di esempio: illustra il metodo main
   e i passi di sviluppo di un'applicazione Java. */
public class HelloWorld {

    // metodi della classe
    public static void main(String[] args) {
        System.out.println("Welcome to Java!");
    }
}
\end{lstlisting}

\vspace{0.2cm}
Un'\textbf{applicazione} Java è un programma eseguibile autonomamente. Ha bisogno di:
\begin{itemize}\itemsep=0pt
    \item almeno una \textbf{classe} (qui \texttt{HelloWorld}) dichiarata dal programmatore\ldots
    \item \ldots che contenga un metodo \textbf{\texttt{main}} con \emph{esattamente} questa firma: è il punto di avvio.
\end{itemize}
\end{frame}

% ========================================================================
\begin{frame}[fragile]{Anatomia riga per riga}
\begin{columns}[T]
\column{0.48\textwidth}
\begin{lstlisting}
public class HelloWorld {
\end{lstlisting}
\small
\begin{itemize}\itemsep=0pt
    \item \texttt{class} introduce la dichiarazione di classe, seguita dal \textbf{nome}
    \item \texttt{public} è un \emph{modificatore di accesso} (ne parleremo)
    \item le graffe \texttt{\{ \}} delimitano il \textbf{corpo} della classe
\end{itemize}

\vspace{0.2cm}
\normalsize
\begin{lstlisting}
public static void main(String[] args)
\end{lstlisting}
\small
\begin{itemize}\itemsep=0pt
    \item \texttt{public}: invocabile dall'esterno (dalla JVM)
    \item \texttt{static}: invocabile \textbf{senza creare oggetti}
    \item \texttt{void}: non restituisce nulla
    \item \texttt{String[] args}: gli argomenti da riga di comando
\end{itemize}

\column{0.48\textwidth}
\normalsize
\begin{lstlisting}
System.out.println("Welcome to Java!");
\end{lstlisting}
\small
\begin{itemize}\itemsep=0pt
    \item \texttt{System.out} è l'oggetto \emph{standard output}
    \item \texttt{println} è un suo \textbf{metodo}: stampa l'argomento e va a capo (\texttt{print} non va a capo)
    \item \texttt{"Welcome to Java!"} è una \textbf{stringa letterale}
    \item ogni \emph{istruzione} termina con \texttt{;}
\end{itemize}

\vspace{0.2cm}
\normalsize
\begin{lstlisting}
/* commento su più righe */
// commento fino a fine riga
\end{lstlisting}
\small I commenti sono ignorati dal compilatore. Usateli per spiegare il \emph{perché}, non il \emph{cosa}.
\end{columns}
\end{frame}

% ========================================================================
\begin{frame}{Regole e convenzioni sui nomi}
\begin{columns}[T]
\column{0.5\textwidth}
\textbf{Regole (le impone il compilatore)}
\begin{itemize}\itemsep=2pt
    \item un identificatore è formato da lettere, cifre, \texttt{\_} e \texttt{\$}
    \item non può iniziare con una cifra né contenere spazi
    \item non può essere una parola chiave (\texttt{class}, \texttt{public}, \texttt{int}\ldots)
    \item Java è \textbf{case sensitive}: \texttt{a1} e \texttt{A1} sono identificatori diversi
    \item il file \texttt{X.java} deve contenere la classe pubblica \texttt{X} (stesso nome, stesse maiuscole)
\end{itemize}

\column{0.5\textwidth}
\textbf{Convenzioni (le impone la comunità)}
\begin{itemize}\itemsep=2pt
    \item classi: \texttt{UpperCamelCase} \textrightarrow{} \texttt{ContoBancario}
    \item metodi e variabili: \texttt{lowerCamelCase} \textrightarrow{} \texttt{calcolaMedia}, \texttt{numeroStudenti}
    \item costanti: \texttt{UPPER\_SNAKE\_CASE} \textrightarrow{} \texttt{MAX\_TENTATIVI}
    \item package: tutto minuscolo \textrightarrow{} \texttt{java.util}
\end{itemize}

\vspace{0.3cm}
\begin{alertblock}{Perché rispettarle}
Il compilatore non se ne accorge, ma chi legge il vostro codice (incluso il voi del futuro) sì. All'esame le convenzioni fanno parte della valutazione.
\end{alertblock}
\end{columns}
\end{frame}

% ========================================================================
\begin{frame}[fragile]{Dal sorgente all'output: dal vivo, da terminale}
\begin{lstlisting}[style=shell]
$ ls
HelloWorld.java
$ javac HelloWorld.java          # compilazione: produce il bytecode
$ ls
HelloWorld.class  HelloWorld.java
$ java HelloWorld                # esecuzione: nome della CLASSE, senza .class
Welcome to Java!
\end{lstlisting}

\begin{mermaid}
flowchart LR
  A["HelloWorld.java<br/><i>sorgente (testo)</i>"] -->|"javac HelloWorld.java"| B["HelloWorld.class<br/><i>bytecode (binario)</i>"]
  B -->|"java HelloWorld"| C["JVM<br/>caricamento · verifica · esecuzione"]
  C --> D[/"Welcome to Java!"/]
  style A fill:#dbeafe,stroke:#1d4ed8
  style B fill:#fef3c7,stroke:#b45309
  style D fill:#dcfce7,stroke:#15803d
\end{mermaid}

\small Il file \texttt{.class} \textbf{non} è un file di testo: aprirlo con un editor mostra caratteri senza senso. Provate.
\end{frame}

% ========================================================================
\begin{frame}[fragile]{Gli errori della prima settimana (e come leggerli)}
\footnotesize
\begin{columns}[T]
\column{0.5\textwidth}
\textbf{1. Nome del file diverso dalla classe}
\begin{lstlisting}[style=shell]
$ javac hello.java
hello.java:1: error: class HelloWorld is
public, should be declared in a file
named HelloWorld.java
\end{lstlisting}

\textbf{2. Manca un \texttt{;} o una graffa}
\begin{lstlisting}[style=shell]
HelloWorld.java:6: error: ';' expected
    System.out.println("Welcome")
                                 ^
\end{lstlisting}
Il compilatore indica \emph{file:riga} e mette \texttt{\^{}} dove si è accorto del problema (non sempre dove sta l'errore!).

\column{0.5\textwidth}
\textbf{3. Maiuscole sbagliate}
\begin{lstlisting}[style=shell]
HelloWorld.java:6: error: package system
does not exist
    system.out.println("Welcome");
\end{lstlisting}

\textbf{4. Eseguire con \texttt{.class} o senza compilare}
\begin{lstlisting}[style=shell]
$ java HelloWorld.class
Error: Could not find or load main class
HelloWorld.class
\end{lstlisting}

\textbf{5. Firma del \texttt{main} sbagliata}
\begin{lstlisting}[style=shell]
Error: Main method not found in class
HelloWorld, please define the main method
as: public static void main(String[] args)
\end{lstlisting}
\end{columns}

\vspace{0.2cm}
\normalsize
\begin{block}{Regola d'oro}
Leggete \textbf{il primo} messaggio d'errore, non l'ultimo: gli errori successivi sono spesso conseguenza del primo.
\end{block}
\end{frame}

% ========================================================================
\section{Esercitazione guidata: programmi interattivi}

\begin{frame}[fragile]{Obiettivo dell'esercitazione}
Finora il programma \emph{parla} soltanto (\texttt{println}). Ora vogliamo programmi che \textbf{dialogano} con l'utente:
\begin{enumerate}\itemsep=0pt
    \item \textbf{leggono} numeri e testo digitati da tastiera;
    \item li \textbf{elaborano} in metodi separati, ognuno con un compito preciso;
    \item \textbf{stampano} il risultato ben formattato.
\end{enumerate}

\MermaidAdjustBoxOpts{max width=0.9\linewidth,max height=0.34\textheight,center}
\begin{mermaid}
flowchart LR
  U(["Utente<br/>digita 7 e 5"]) --> K["Tastiera<br/>System.in"]
  K -->|"nextInt()"| S["<b>Scanner</b><br/>Byte → int"]
  S -->|"7, 5"| M["<b>Metodo somma(a, b)</b><br/>Calcola 12"]
  M -->|"12"| P["<b>System.out.printf</b><br/>Formatta"]
  P --> O(["Schermo<br/>La somma è: 12"])
  style S fill:#dbeafe,stroke:#1d4ed8
  style M fill:#fef3c7,stroke:#b45309
  style P fill:#dcfce7,stroke:#15803d
\end{mermaid}
\MermaidAdjustBoxOpts{max width=0.95\linewidth,max height=0.62\textheight,center}

\vspace{0.1cm}
Servono tre ingredienti nuovi, che vediamo uno alla volta e poi mettiamo insieme in \texttt{SommaNumeri}:
\begin{itemize}\itemsep=0pt
    \item \texttt{System.out.printf} -- stampa formattata (colonne, decimali)
    \item \texttt{Scanner} -- lettura da tastiera
    \item metodi \texttt{static} -- ``funzioni'' per spezzare il programma in compiti
\end{itemize}
\end{frame}

% ========================================================================
\begin{frame}[fragile]{Ingrediente 1 -- Stampa formattata con \texttt{printf}}
Si comporta come la \texttt{printf} del C:
\begin{lstlisting}
int n = 7;
double media = 27.3333;
String nome = "Anna";
System.out.printf("Ciao %s, hai %d esami con media %.2f%n", nome, n, media);
// Ciao Anna, hai 7 esami con media 27.33
\end{lstlisting}

\begin{center}\small
\begin{tabular}{ll}
\toprule
\textbf{Specificatore} & \textbf{Significato} \\
\midrule
\texttt{\%d} & intero \\
\texttt{\%f}, \texttt{\%.2f} & reale (con 2 decimali) \\
\texttt{\%s} & stringa \\
\texttt{\%c} & carattere \\
\texttt{\%n} & a capo (portabile: usatelo al posto di \texttt{\textbackslash n}) \\
\texttt{\%5d}, \texttt{\%-10s} & larghezza minima 5, allineato a sinistra su 10 \\
\bottomrule
\end{tabular}
\end{center}

\small In alternativa: \texttt{System.out.println("Ciao " + nome + ", media " + media);} -- l'operatore \texttt{+} concatena stringhe (ne parleremo la prossima lezione).
\end{frame}

% ========================================================================
\begin{frame}[fragile]{Ingrediente 2 -- Leggere dall'utente con \texttt{Scanner}}
\begin{lstlisting}
import java.util.Scanner;   // 1. la classe Scanner sta nel package java.util

public class Eta {
    public static void main(String[] args) {
        Scanner in = new Scanner(System.in);   // 2. un oggetto Scanner "collegato" alla tastiera
        System.out.print("Come ti chiami? ");
        String nome = in.nextLine();           // 3. legge un'intera riga
        System.out.print("Quanti anni hai? ");
        int eta = in.nextInt();                // 4. legge il prossimo intero
        System.out.printf("Ciao %s, tra un anno avrai %d anni.%n", nome, eta + 1);
        in.close();                            // 5. chiusura (buona pratica)
    }
}
\end{lstlisting}

\begin{columns}[T]
\column{0.5\textwidth}
\small
\begin{itemize}\itemsep=0pt
    \item \texttt{import}: dice al compilatore dove trovare \texttt{Scanner}
    \item \texttt{new}: \textbf{crea un oggetto} (lo studieremo a fondo); \texttt{System.in} è lo standard input
\end{itemize}
\column{0.5\textwidth}
\small
\textbf{Metodi principali}: \texttt{nextInt()}, \texttt{nextDouble()}, \texttt{next()} (una parola), \texttt{nextLine()} (una riga), \texttt{hasNextInt()} (c'è un intero da leggere?)
\end{columns}
\end{frame}

% ========================================================================
\begin{frame}[fragile]{Ingrediente 3 -- Metodi \texttt{static}: funzioni per organizzare il codice}
\begin{columns}[T]
\column{0.58\textwidth}
\begin{lstlisting}[basicstyle=\ttfamily\scriptsize]
public class Geometria {

    // modificatori, tipo, nome, parametri
    public static double areaCerchio(double r) {
        return Math.PI * r * r;
    }

    public static void stampaSaluto() {  // void
        System.out.println("Ciao!");
    }

    public static void main(String[] args) {
        stampaSaluto();               // chiamata
        double a = areaCerchio(2.0);  // 12.566...
        System.out.printf("Area: %.3f%n", a);
    }
}
\end{lstlisting}

\column{0.42\textwidth}
\small
\textbf{\texttt{static}} significa che il metodo \textbf{appartiene alla classe}, non a un oggetto: si può chiamare senza fare \texttt{new}.

\vspace{0.2cm}
Per ora \emph{tutti} i nostri metodi saranno \texttt{static}: sono le ``funzioni'' che conoscete dal C. Dalla lezione~5 scopriremo i metodi \emph{di istanza}.

\vspace{0.2cm}
\begin{block}{Perché spezzare il codice in metodi}
\begin{itemize}\itemsep=0pt
    \item un metodo = \textbf{un compito}, con un nome che lo descrive
    \item si \emph{riusa} e si \emph{testa} separatamente
    \item il \texttt{main} diventa la ``trama'', leggibile in 5 righe
\end{itemize}
\end{block}
\end{columns}
\end{frame}

% ========================================================================
\begin{frame}[fragile]{Commentiamo insieme -- \texttt{SommaNumeri}}
\begin{columns}[T]
\column{0.6\textwidth}
\begin{lstlisting}[basicstyle=\ttfamily\scriptsize]
import java.util.Scanner;

public class SommaNumeri {

    public static void main(String[] args) {
        int risultato = sommaDueNumeri();
        System.out.printf("La somma è: %d%n", risultato);
    }

    public static int sommaDueNumeri() {
        Scanner scanner = new Scanner(System.in);
        System.out.printf("Primo numero: ");
        int numero1 = scanner.nextInt();
        System.out.printf("Secondo numero: ");
        int numero2 = scanner.nextInt();
        scanner.close();
        return numero1 + numero2;
    }
}
\end{lstlisting}

\column{0.4\textwidth}
\small
\textbf{Leggiamolo dall'alto}
\begin{enumerate}\itemsep=2pt
    \item il \texttt{main} \emph{chiama} \texttt{sommaDueNumeri()} e mette il valore restituito in \texttt{risultato}
    \item \texttt{sommaDueNumeri} è \texttt{static} e restituisce un \texttt{int}
    \item lo \texttt{Scanner} traduce i byte digitati in un \texttt{int} utilizzabile
    \item \texttt{return} termina il metodo e ``consegna'' il valore al chiamante
\end{enumerate}

\textbf{Poi lo miglioriamo insieme}
\begin{itemize}\itemsep=2pt
    \item separare \emph{input} e \emph{calcolo}: \texttt{static int somma(int a, int b)}
    \item leggere i numeri nel \texttt{main} e lasciare ai metodi solo il calcolo
    \item cosa succede se l'utente scrive \texttt{ciao}? (lo risolveremo alla lezione~10)
\end{itemize}
\end{columns}
\end{frame}

% ========================================================================
\begin{frame}[fragile]{Variante -- Argomenti da riga di comando}
Il parametro \texttt{args} del \texttt{main} contiene ciò che scriviamo dopo il nome della classe:
\begin{lstlisting}
public class Saluta {
    public static void main(String[] args) {
        if (args.length == 0) {
            System.out.println("Uso: java Saluta <nome> [<nome> ...]");
            return;
        }
        for (int i = 0; i < args.length; i++) {   // args.length: quanti argomenti
            System.out.printf("Ciao, %s!%n", args[i]);
        }
    }
}
\end{lstlisting}
\begin{lstlisting}[style=shell]
$ java Saluta Anna Luca
Ciao, Anna!
Ciao, Luca!
$ java Saluta
Uso: java Saluta <nome> [<nome> ...]
\end{lstlisting}
\end{frame}

% ========================================================================
\begin{frame}[fragile]{Insieme, passo passo (25 min) -- \texttt{Media.java}}
\textbf{Il problema.} Chiedere quanti voti si vogliono inserire, leggerli uno alla volta e stampare la media con 2 decimali.

\vspace{0.3cm}
\textbf{Prima di scrivere codice, decidiamo insieme:}
\begin{enumerate}\itemsep=3pt
    \item quali \textbf{dati} servono (quanti voti? la somma? di che tipo?);
    \item quale \textbf{calcolo} isolare in un metodo, e cosa deve ricevere e restituire\\
          \small \textrightarrow{} \texttt{static double media(double somma, int n)}\normalsize;
    \item cosa fa il \texttt{main}: crea lo \texttt{Scanner}, legge, chiama il metodo, stampa.
\end{enumerate}

\vspace{0.3cm}
\begin{block}{Come procediamo}
Lo scriviamo \textbf{insieme}, una riga alla volta: io al PC proiettato, voi sul vostro portatile in VS Code. Alla fine compiliamo ed eseguiamo (almeno una volta \textbf{da terminale}) e proviamo qualche input strano. Quel che non finiamo in aula, lo finite a casa.
\end{block}
\end{frame}

% ========================================================================
\begin{frame}[fragile]{Una possibile soluzione -- \texttt{Media.java}}
\begin{lstlisting}[basicstyle=\ttfamily\scriptsize]
import java.util.Scanner;

public class Media {

    // il calcolo, separato dall'input/output
    public static double media(double somma, int n) {
        return somma / n;   // somma e` double: la divisione e` reale
    }

    public static void main(String[] args) {
        Scanner in = new Scanner(System.in);
        System.out.print("Quanti voti? ");
        int n = in.nextInt();
        double somma = 0;
        for (int i = 1; i <= n; i++) {
            System.out.print("Voto " + i + ": ");
            somma += in.nextInt();
        }
        System.out.printf("Media: %.2f%n", media(somma, n));
        in.close();
    }
}
\end{lstlisting}
\vspace{-0.2cm}\small Perché \texttt{somma} è \texttt{double} e non \texttt{int}? Provate a cambiarlo e a inserire i voti 27 e 28. E cosa succede con $n = 0$?
\end{frame}

% ========================================================================
% ========================================================================
\section{Wrap-up}

\begin{frame}{Riepilogo: cosa abbiamo trattato}
\small
\begin{center}
\renewcommand{\arraystretch}{1.3}
\begin{tabular}{|l|l|}
\hline
\textbf{Concetto (unità di riferimento)} & \textbf{Dove compare oggi} \\
\hline
Bytecode, JVM, WORA (U2) & diagrammi C++ vs Java; le 5 fasi; dentro la JVM \\
\hline
JDK vs JRE, versioni LTS (U1, U2) & verifica dell'installazione, \texttt{java -version} \\
\hline
Classe, \texttt{main}, \texttt{System.out} (U1) & anatomia di \texttt{HelloWorld} \\
\hline
\texttt{javac} / \texttt{java}, file \texttt{.class} (U1) & compilazione ed esecuzione da terminale \\
\hline
Identificatori, case sensitivity, CamelCase (U1) & regole e convenzioni; errori tipici \\
\hline
\texttt{import}, \texttt{Scanner}, \texttt{printf} (U6) & \texttt{Eta}, \texttt{SommaNumeri}, \texttt{Media} \\
\hline
Metodi \texttt{static}, \texttt{return} (U6) & \texttt{Geometria}, \texttt{sommaDueNumeri}, \texttt{Media} \\
\hline
\texttt{args}, \texttt{args.length} (U1) & \texttt{Saluta} \\
\hline
\end{tabular}
\end{center}

\vspace{0.3cm}
\textbf{Prossima lezione (22 settembre)}: il sistema dei tipi -- primitivi, costanti, riferimenti e \texttt{null} -- e le stringhe (\texttt{String}, \texttt{StringBuilder}).
\end{frame}

% ========================================================================
\begin{frame}{Homework 1 (da svolgere entro lunedì 21 settembre)}
\begin{enumerate}
    \item \textbf{\texttt{Calcolatrice.java}}: legge due numeri reali e un operatore (\texttt{+ - * /}) e stampa il risultato con 3 decimali. Un metodo statico per ogni operazione; gestite la divisione per zero con un messaggio.
    \item \textbf{\texttt{Bmi.java}}: legge peso (kg) e altezza (m), calcola l'indice di massa corporea $\text{BMI} = \text{peso} / \text{altezza}^2$ in un metodo \texttt{static double bmi(double peso, double altezza)} e stampa la categoria: \emph{sottopeso} se BMI $< 18{,}5$; \emph{normopeso} tra $18{,}5$ e $24{,}9$; \emph{sovrappeso} tra $25$ e $29{,}9$; \emph{obesità} se BMI $\geq 30$.
    \item \textbf{\texttt{Fizzbuzz.java}}: legge $n$ da \texttt{args} (default 30) e stampa i numeri da 1 a $n$ sostituendo i multipli di 3 con \texttt{Fizz}, di 5 con \texttt{Buzz}, di entrambi con \texttt{FizzBuzz}.
\end{enumerate}

\vspace{0.2cm}
\begin{block}{Come funziona}
Non è obbligatorio, ma fortemente consigliato: martedì prossimo chiederò a qualcuno di aprire il proprio codice e presentarlo alla classe, e lo discutiamo insieme. Tenete i file \texttt{.java} in una cartella \texttt{hw1/} del progetto VS Code; ogni file deve compilare ed eseguire. Nomi di classi, metodi e variabili secondo le convenzioni viste oggi.
\end{block}
\end{frame}

% ========================================================================
\begin{frame}
\begin{center}
{\Huge Domande?}\\[1.5em]
{\Large Buon lavoro!}\\[2em]
{\normalsize \href{mailto:mattia.fonisto@unina.it}{\texttt{mattia.fonisto@unina.it}}}
\end{center}
\end{frame}

\end{document}
